% Source: https://artofproblemsolving.com/community/c3h2180650_problem_set_revivals
% Stable identity: aops-topic-2180650-q11
\textbf{With the formula.} $\dbinom nk=\dfrac{n!}{k!\,(n-k)!}$, and replacing $k$ by $n-k$ gives $\dbinom n{n-k}=\dfrac{n!}{(n-k)!\,\left(n-(n-k)\right)!}=\dfrac{n!}{(n-k)!\,k!}$, the same number.

\textbf{By counting.} $\dbinom nk$ counts the ways to choose $k$ of $n$ people to be on a committee. Choosing the $k$ people who are on the committee is the same as choosing the $n-k$ people who are left out: each choice of one determines the other (a subset corresponds to its complement). So the number of ways to choose $k$ people equals the number of ways to choose $n-k$ people, which is $\dbinom n{n-k}$.
